\section{The scalar gauge inverse and bounded current vertices}
\label{gauge-section}

The fine inverse alone does not bound the current perturbation.
The source gauge contains derivatives, so we first control its
scalar constrained inverse in the same physical cell norms.
Keep the matched scalar average $Q_0$, scalar Laplacian
$\mathcal L=D^*D$, and the source operators
\[
 G'=(\mathcal L+Q_0^*a'Q_0)^{-1},\qquad
 R=I-G'Q_0^*(Q_0G'^2Q_0^*)^{-1}Q_0G',
 \qquad Z=G'RG'.
\]
All adjoints use the physical measures. Complex continuations use
the fixed bilinear transpose, never conjugation of a varying parameter.

\subsection{The constrained biharmonic inverse}

The operator
\[
 Z=G'^2-G'^2Q_0^*(Q_0G'^2Q_0^*)^{-1}Q_0G'^2
\]
is the inverse of the quadratic form $\norm{\mathcal L\lambda}^2$
on $\ker Q_0$, extended by zero on its orthogonal complement.
Indeed it is the usual constrained inverse of $G'^{-2}$, and for
$Q_0\lambda=0$ one has
$\norm{G'^{-1}\lambda}^2=\norm{\mathcal L\lambda}^2$.
It follows that $Q_0Z=0=ZQ_0^*$ and that $Z$ is independent of
the auxiliary positive penalty $a'$. The constant scalar mode is
excluded by the constraint, so no inverse is assigned to that mode.

On a real nonzero Bloch fiber the constrained inverse is
\[
 Z_{rs}=\frac{\delta_{rs}}{P_r^2}
 -\frac{a_r^-a_s}{P_r^2P_s^2W_2},
 \qquad W_2=\frac{D_{s,2}}{P_0^2}.
\]
Its completely canceled form is
\begin{align}
 Z_{00}&=\frac{R_{s,2}}{D_{s,2}},
 &Z_{0s}&=-\frac{a_0^-a_s}{P_s^2D_{s,2}}\quad(s\ne0),
 \label{gauge-Z-central}\\
 Z_{r0}&=-\frac{a_r^-a_0}{P_r^2D_{s,2}}\quad(r\ne0),
 &Z_{rs}&=\frac{\delta_{rs}}{P_r^2}
       -\frac{P_0^2a_r^-a_s}{P_r^2P_s^2D_{s,2}}
          \quad(r,s\ne0).
 \label{gauge-Z-high}
\end{align}
The notation $D_{s,2}$ denotes the scalar denominator, not a
dependence on the second matrix index $s$. No central inverse
power of $P_0$ remains. At zero coarse momentum the high averaging
coefficients vanish: $Z$ is zero on the constant mode and equals
$P_r^{-2}$ on the nonzero aliases. For $M=1$, it is identically zero.

\begin{proposition}\label{gauge-Z-locality}
Every placement of at most four total fine forward or backward
differences on the two sides of $Z$ has a uniformly bounded
holomorphic physical-cell symbol and a uniformly exponentially
decaying cell-operator kernel. The bounds are independent of
refinement and compatible periodic volume.
\end{proposition}

\begin{proof}
Let the output order be $j$ and the input order be $l$, with
$j+l\le4$. The high diagonal in \eqref{gauge-Z-high} is bounded
by $C|r|^{j+l-4}$. Its high rank vectors are bounded by
$C|a_r^-||r|^{j-4}$ and $C|a_r||r|^{l-4}$.
Each order is at most four, so alias Parseval bounds both vector
Hilbert norms uniformly. The central row and column use the same
vectors; central multipliers are bounded. The scalar denominator
$D_{s,2}$ stays uniformly away from zero on the common strip.
For a complex rank-one matrix $vw^T$, its Hilbert norm is bounded
by $\norm v_2\norm w_2$, without conjugating the analytic symbol.

These are operator bounds before reconstruction and hence do not
count the growing number of aliases. The finite Fourier reconstruction
and its inverse are uniformly bounded in physical cell Hilbert
norm, with the inverse containing its required $M^{-d}$ factor.
Alias permutation and the integer-site phases give periodic gluing
on the enlarged charts. Contour shifting gives exponential cell
decay; summing at a smaller exponent gives weighted Schur bounds.
Fourier sampling gives exact periodizations on finite tori, including
the stated zero mode and small-period incidence multiplicities.
\end{proof}

The source definitions give
\[
 Q_0G'R=0,\quad \mathcal L G'R=R,\quad
 RG'\mathcal L=R.
\]
Consequently
\[
 \Lambda:=G'RD^*=Z\mathcal L D^*,\qquad
 D\Lambda=DZ\mathcal L D^*.
\]
The total fine derivative orders are three and four, respectively.
Proposition~\ref{gauge-Z-locality} therefore bounds both operators
in a common positive weighted cell-Schur norm. This is a Hilbert
operator conclusion; it does not assert pointwise bounds for four
fine derivatives.

\subsection{A current perturbation with no derivative of the current}

At the flat reference the deposited gauge BCH coefficients are
\[
 c^{(1)}(A,\lambda)(b)
 =\frac i2[\lambda(b_-)+\lambda(b_+),A(b)],
\]
\[
 2c^{(2)}(\lambda_1,\lambda_2)(b)
 =\frac1{2i}\bigl([\lambda_1(b_-),D\lambda_2(b)]
                 +[\lambda_2(b_-),D\lambda_1(b)]\bigr).
\]
Define $\mathfrak j_J$, $\mathfrak k_J=\mathfrak j_J^*$ and
the symmetric scalar form $\mathfrak m_J$ by pairing these two
expressions with the fine current $J$. More explicitly, define
\[
 \langle\lambda,\mathfrak j_J A\rangle
   =\langle c^{(1)}(A,\lambda),J\rangle,\qquad
 \langle\lambda_1,\mathfrak m_J\lambda_2\rangle
   =2\langle c^{(2)}(\lambda_1,\lambda_2),J\rangle.
\]
Also define
the finite-incidence map $V_J$ through
\[
 \langle\lambda,V_J A\rangle
 =\frac1{2i}\sum_b\eta^d
           \langle[\lambda(b_-),A(b)],J(b)\rangle.
\]
Then, as a quadratic form,
\[
 \mathfrak m_J=V_JD+D^*V_J^*.
\]
The fixed Lie bracket norm and finite incidence give
\[
 \norm{\mathfrak j_J}_{2\to2}+\norm{V_J}_{2\to2}
 \le C_d\norm J_\infty.
\]
For example, there are at most $2d$ incident contributions at a
site, and each bond has two endpoints. Summing with the common
fine volume measure pays exactly those fixed multiplicities.
There is no power of $M$. The maps have a fixed cell range, so
the same bound holds in each fixed weighted cell-Schur norm.

The auxiliary current form of the deposited source is
\begin{equation}\label{gauge-Theta}
 \begin{split}
 \Theta'_J
 &=\mathfrak k_J\Lambda+\Lambda^*\mathfrak j_J
                    +\Lambda^*\mathfrak m_J\Lambda\\
 &=\mathfrak j_J^*\Lambda+\Lambda^*\mathfrak j_J
   +\Lambda^*V_JD\Lambda+(D\Lambda)^*V_J^*\Lambda.
 \end{split}
\end{equation}
Submultiplicativity of the weighted algebra proves, for a fixed
positive exponent $\kappa$,
\begin{equation}\label{gauge-Theta-bound}
 \norm{\Theta'_J}_\kappa\le\theta\norm J_\infty,
\end{equation}
uniformly in $M$ and volume. The expression is linear and
holomorphic in complex $J$. Fine derivatives act on the scalar
legs through $D\Lambda$; no derivative of $J$ has been taken.
If $J$ is supported in one coarse cell, its coefficient maps read
that cell and write only fixed neighboring cells. This marked
support will give the localized current derivatives below.

\subsection{The zero-current cancellation}

Let $E_0=G_0Q^*q_0$ be the normalized source extension, with
$q_0=q_M+aI$. Its source gauge gives
$\Lambda E_0=0$. Hence
\[
 \Theta'_J E_0=\Lambda^*\mathfrak j_J E_0,
 \qquad E_0^*\Theta'_J E_0=0.
\]
For $\widetilde G_0=G_0-E_0q_0^{-1}E_0^*$ the identities
\begin{equation}\label{gauge-response-identities}
 G_0\Lambda^*=\widetilde G_0\Lambda^*=DZ,
 \qquad \Lambda DZ=Z
\end{equation}
follow directly: for $\lambda=Zf$,
$Q_0\lambda=0$ and $\mathcal L\lambda=RG'f$, so
\[
 (K+DRD^*+Q^*aQ)D\lambda=DRG'f=\Lambda^*f,
 \qquad QD\lambda=0.
\]
In particular, exact differentiation of the auxiliary family at
zero current gives
\[
 Dq_0[J]=0,\qquad DE_0[J]=DZ\mathfrak j_JE_0,
\]
\begin{equation}\label{gauge-current-hessian}
 \begin{split}
 D^2q_0[J_1,J_2]
 ={}&-(\mathfrak j_{J_1}E_0)^*Z(\mathfrak j_{J_2}E_0)\\
    &-(\mathfrak j_{J_2}E_0)^*Z(\mathfrak j_{J_1}E_0).
 \end{split}
\end{equation}
Here $Dq_0$ means the current derivative of the auxiliary $q_J$
at $J=0$, rather than a derivative of the fixed penalty. The first
extension response is a fine gradient with zero coarse average and
zero curl. The quadratic precision response can be nonzero.
The separate constraint-current term in the full source family
is not canceled by these identities.
